\documentclass{article}

\usepackage[preprint]{neurips_2024}
\usepackage[utf8]{inputenc}
\usepackage[T1]{fontenc}
\usepackage{hyperref}
\usepackage{url}
\usepackage{booktabs}
\usepackage{amsfonts}
\usepackage{amsmath}
\usepackage{amssymb}
\usepackage{amsthm}
\usepackage{mathtools}
\usepackage{bm}
\usepackage{nicefrac}
\usepackage{microtype}
\usepackage{xcolor}
\usepackage{graphicx}
\usepackage{algorithm}
\usepackage{algpseudocode}
\usepackage{float}
\usepackage{multirow}
\usepackage{array}
\usepackage{makecell}
\usepackage{tabularx}
\hypersetup{hypertexnames=false}

\newcommand{\R}{\mathbb{R}}
\newcommand{\method}{\textsc{H2S-NN}}
\newcommand{\ampss}{\textsc{AMPSS}}
\newcommand{\thspool}{\textsc{THS-Pool}}
\newcommand{\cF}{\mathcal{F}}
\newcommand{\cL}{\mathcal{L}}
\newcommand{\norm}[1]{\left\lVert #1\right\rVert}
\newcommand{\transpose}{\mathsf{T}}
\newcommand{\concat}{\mathbin\Vert}
\newtheorem{proposition}{Proposition}

\title{H2S-NN: Heterogeneous Sheaf--Hodge Neural Networks\\
with Adaptive Schema-to-Cell Selection}

% Replace the affiliation placeholder before submission.
\author{
  Peeradon Sarnkaew \\
  \texttt{peeradonsarnkaew46@gmail.com}
}

\begin{document}
\maketitle

\begin{abstract}
Heterogeneous graphs combine several node and relation types, while many of their observed
relations are heterophilic. Conventional heterogeneous graph neural networks represent typed
paths as message-passing channels, and existing heterogeneous sheaf networks learn a geometry on
the observed graph but remain confined to its nodes and edges. We introduce \method, a
heterogeneous sheaf--Hodge architecture in which concrete typed path instances provide
higher-order, face-like supports for diffusion over an explicit edge cochain. Relation-conditioned
node-to-edge maps define a learned first coboundary, an ordered sequence encoder constructs
edge-to-face maps, and matrix-free lower and upper Hodge terms evolve node and edge states without
materializing global block matrices. We further introduce \ampss, a two-stage selection procedure
that enumerates typed schemas, ranks them using training-only task signals and structural coverage,
and refines the retained schemas and instances through differentiable nonnegative gates inside the
upper edge Laplacian. We state precisely where the current construction differs from a strict
cellular cochain complex: synthetic closure pairs correct the node Laplacian, whereas the learned
second coboundary remains path-edge-only and $\delta_1\delta_0=0$ is encouraged rather than
guaranteed. Across the completed evaluation, \method{} obtains the strongest reported average in
all six node-classification columns, all eight node-clustering columns, and all eight
link-prediction columns; it ties the MUTAG ceiling and leads the remaining graph-classification
comparisons. On ACM it also records the lowest runtime, parameter count, and peak memory among the
compared systems. A DBLP ablation attributes additional gains to automatic schema selection and
instance-level refinement. The in-development \thspool{} hierarchical graph-pooling extension is
intentionally reserved as a placeholder and is not used to support these claims.
\end{abstract}

\section{Introduction}
\label{sec:introduction}

Heterogeneous graphs are a natural representation in many real-world relational data, such as citation networks, or social networks. 
In heterogeneous graphs, they consist of multiple types of nodes and edges, which can represent different entities and relationships.

The graph neural networks (GNNs) have been widely used for various graph-based tasks, such as node classification, link prediction, and graph classification.
On those datasets, Treating such a network as homogeneous erases this
distinction. Treating every observed edge as an instruction to average neighbouring features
creates a second difficulty: heterogeneous relations are often heterophilic, so adjacent entities
need not share a label or even inhabit comparable raw feature spaces.

However, with recent research, the study found that they also exist a property called heterophily, which means that connected nodes are often dissimilar in their features or labels.
Therefore heterophilic heterogeneous graphs present a challenge for traditional graph neural networks (GNNs), which often assume that connected nodes are similar and can be aggregated together.


Two established approaches address complementary parts of this problem. Heterogeneous GNNs use
typed parameters, attention, or meta-paths to preserve semantic information
\citep{wang2019han,fu2020magnn,hu2020hgt}. Meta-path models are particularly attractive because a
typed path such as paper--author--paper describes a relation that is not visible from either edge in
isolation. Nevertheless, the path usually serves only as a route for collecting neighbours; after
aggregation, the concrete path instance has no further mathematical role. Sheaf neural networks
instead attach local vector spaces and learned restriction maps to a graph, replacing direct
feature averaging with comparison after learned transport. Neural Sheaf Diffusion demonstrated
the value of this geometry for heterophily and oversmoothing \citep{bodnar2022neural}, and HetSheaf
specialized restriction-map prediction to heterogeneous graphs \citep{braithwaite2024hetsheaf}.
These models, however, operate on the observed 1-skeleton.

\method{} connects the two perspectives by changing what a typed path becomes. Each selected
concrete path instance is retained as a face-like higher-order support. Its ordered boundary edges
are encoded into restriction maps, enter a learned second coboundary $\delta_1$, and therefore
modify the upper term $\delta_1^{\transpose}\delta_1$ of an edge Hodge operator. The model maintains
both a node cochain and an independently initialized edge cochain. It diffuses the former with a
learned $L_0$ and the latter with both lower and upper components of $L_1$. Thus the selected path
affects the operator itself rather than merely adding another neighbourhood view.

Manual meta-paths are useful when reliable domain knowledge exists, but they are difficult to
scale across schemas. We therefore develop \ampss, an automatic candidate-pool construction and
refinement procedure. A bounded type-graph search first proposes schemas. Training-only task
signals, coverage, diversity, and redundancy determine a compact pool. Finally, a schema gate
$\alpha_\sigma$ and an instance gate $\beta_f$ weight upper-Laplacian summands during training. The
combinatorial preselection is not differentiable, while refinement of the fixed pool is
differentiable through the task and regularization objectives.

The contributions of this work are as follows.
\begin{enumerate}
  \item We formulate typed path instances as operator-level higher-order supports for a
  heterogeneous sheaf model. Learned edge-to-face maps make the selected paths act through the
  upper edge Laplacian rather than through a semantic-channel sum.
  \item We introduce a matrix-free heterogeneous sheaf--Hodge backbone that learns
  relation-conditioned node-to-edge maps, initializes an explicit edge cochain, and performs
  diffusion over node and edge ranks.
  \item We propose \ampss, a two-stage schema-to-cell selection mechanism with training-only
  preselection, hierarchical schema/instance gates, and partial-schema propagation.
  \item We give a code-grounded account of the construction's guarantees and limitations,
  including positive semidefiniteness of the weighted Hodge terms, the absence of an exact chain
  condition, and the task-dependent visibility of the edge cochain.
  \item We evaluate the full discovery-enabled model on node classification, graph classification,
  node clustering, and heterogeneous link prediction, together with a schema-selection ablation
  and an ACM resource comparison.
\end{enumerate}

Automatic meta-path generation alone is not the novelty: LMSPS and differentiable meta-graph or
meta-multigraph search already automate schema selection \citep{yan2024lmsps,jiang2022pmmm}. Our
distinction is operator-level: selected instances become inputs to a learned second coboundary and
thereby reshape edge diffusion.

\section{Background and Related Work}
\label{sec:background}

\subsection{Heterogeneous graphs and typed paths}

A heterogeneous graph is a tuple $G=(V,E,\phi,\psi)$, where $\phi:V\rightarrow\mathcal{T}$ and
$\psi:E\rightarrow\mathcal{R}$ assign node and relation types. A meta-path schema
$\sigma=(t_0,r_1,t_1,\ldots,r_\ell,t_\ell)$ specifies a valid walk in the type graph. A concrete
instance $f=(v_0,e_1,v_1,\ldots,e_\ell,v_\ell)$ follows that schema in $G$. HAN and MAGNN aggregate
features over manually specified meta-path views \citep{wang2019han,fu2020magnn}, whereas HGT uses
typed attention without requiring explicit paths \citep{hu2020hgt}. LMSPS progressively searches
long-range paths, and PMMM searches a differentiable meta-multigraph while using partial message
propagation to improve stability \citep{yan2024lmsps,jiang2022pmmm}. These methods search or weight
communication structures; they do not retain each path instance as an element of a second
coboundary operator.

\subsection{Cellular sheaves and neural diffusion}

A cellular sheaf $\cF$ on a graph assigns a vector space $\cF(v)$ to each node, a vector space
$\cF(e)$ to each edge, and a linear restriction map $F_{v\unlhd e}:\cF(v)\rightarrow\cF(e)$ for
each incident pair. With a chosen edge orientation, the first coboundary compares transported
endpoint signals,
\begin{equation}
  (\delta_0 x)_e = F_{t(e)\unlhd e}x_{t(e)}-F_{s(e)\unlhd e}x_{s(e)}.
  \label{eq:background-delta0}
\end{equation}
The associated sheaf Laplacian is $L_0=\delta_0^{\transpose}\delta_0$. When all stalks are
one-dimensional and all restrictions are identities, this reduces to the ordinary graph
Laplacian. Neural Sheaf Diffusion learns the restrictions from data and uses the induced geometry
to control diffusion in heterophilic settings \citep{bodnar2022neural}. HetSheaf conditions the
predictor on node and edge types and develops a gauge-aware graph readout
\citep{braithwaite2024hetsheaf}. \method{} retains typed restriction learning while adding an
explicit edge cochain and a path-supported second operator.

\subsection{Heterophily-aware heterogeneous learning}

Recent methods address heterogeneous heterophily by rewiring, spectral filtering, latent-graph
construction, contrastive objectives, or cross-type homophily measurements
\citep{guo2023hdhgr,lu2024h2sgnn,shen2024latgrl,wang2025hgms,tao2025cthge}. These approaches are
important baselines because they target the same failure mode of naive neighbour averaging. Their
mechanism is nevertheless different from learning maps on node--edge and edge--face incidences.

\subsection{Positioning}

Table~\ref{tab:positioning} summarizes the boundary of the contribution. We use ``face-like''
deliberately: Section~\ref{sec:closure} explains why the implemented operators form a generalized
path-supported cochain system rather than an exact cellular cochain complex.

\begin{table}[t]
  \centering
  \caption{Position of \method{} relative to adjacent model families.}
  \label{tab:positioning}
  \small
  \begin{tabularx}{\textwidth}{lXX}
    \toprule
    \textbf{Family} & \textbf{Primary mechanism} & \textbf{Difference in \method{}} \\
    \midrule
    HAN / MAGNN & Meta-path views route and aggregate messages. & Concrete path instances remain in an edge-to-face operator. \\
    HGT & Typed attention over the observed graph. & Adds an explicit edge cochain and a learned upper edge term. \\
    LMSPS / PMMM & Automatically select schemas or meta-multigraph channels. & Selected structures weight $\delta_1^{\transpose}\delta_1$, not only aggregation. \\
    Neural Sheaf Diffusion & Learned sheaf diffusion on a homogeneous 1-skeleton. & Conditions on heterogeneous types and introduces path-supported edge diffusion. \\
    HetSheaf & Heterogeneous restriction maps and gauge-aware readout. & Adds ordered edge-to-face maps and an explicit edge Hodge operator. \\
    \bottomrule
  \end{tabularx}
\end{table}

\section{H2S-NN}
\label{sec:method}

\subsection{Problem setting and notation}

Let $N=|V|$ and $M=|E|$. We use a common stalk dimension $d$ for node, edge, and
face-output vectors, a relation-embedding dimension $d_r$, and a sequence hidden dimension $d_h$.
Node and edge cochains are $x\in\R^{N\times d}$ and $s\in\R^{M\times d}$. The $F$ selected path
instances contain $I=\sum_{f=1}^{F}|\partial f|$ path-edge incidences. Superscripts ``src'' and
``tgt'' denote the two node-to-edge restriction maps of an oriented edge;
$A_{e\triangleright f}$ denotes an edge-to-face map. Adjoint operators are written with
$\transpose$ because the implementation uses the Euclidean inner product on every stalk.

\begin{figure}[t]
  \centering
  \setlength{\fboxsep}{6pt}
  \begin{tabular}{@{}c@{\hspace{0.5em}}c@{\hspace{0.5em}}c@{\hspace{0.5em}}c@{}}
    \fbox{\parbox{0.20\textwidth}{\centering\textbf{1. Typed input}\\pad features; encode node and relation types}} &
    \fbox{\parbox{0.20\textwidth}{\centering\textbf{2. Path support}\\manual schemas or \ampss; enumerate instances}} &
    \fbox{\parbox{0.20\textwidth}{\centering\textbf{3. Geometry}\\learn node--edge and edge--face maps}} &
    \fbox{\parbox{0.20\textwidth}{\centering\textbf{4. Prediction}\\Hodge diffusion; task readout}}
  \end{tabular}
  \caption{The \method{} pipeline. Selected path instances act through learned edge-to-face maps
  and the upper component of the edge operator.}
  \label{fig:pipeline}
\end{figure}

\subsection{Typed input representation}

Raw feature matrices may have different widths across node types. The data layer pads them to a
common width and appends relation-level structural statistics derived using training labels only.
For a node $v$ of type $\phi(v)$, the implemented encoder is
\begin{equation}
  h_v^{(0)}=\operatorname{Dropout}\!\left(
  \operatorname{GELU}\!\left(
  \operatorname{LN}\!\left(W_{\phi(v)}x_v+b_{\phi(v)}\right)+q_{\phi(v)}
  \right)\right)\in\R^d,
  \label{eq:node-encoder}
\end{equation}
where each node type has a separate affine map and $q_t$ is a learned type embedding. If a graph
has no node features, a free embedding table supplies one learned $d$-vector per node; this avoids
the bias-only collapse of a zero-width linear layer. Each relation $r$ has an embedding
$z_r\in\R^{d_r}$.

\subsection{Relation-conditioned node-to-edge maps}

For every stored directed edge $e=(u\xrightarrow{r}v)$, two relation-specific MLPs receive
$[h_u^{(0)}\concat h_v^{(0)}\concat z_r]$ and predict maps
$F_e^{\mathrm{src}},F_e^{\mathrm{tgt}}\in\R^{d\times d}$. The default maps are conformal,
\begin{equation}
  F_e=\lambda_e Q_e,\qquad Q_e^{\transpose}Q_e=I,\qquad
  \lambda_e=\exp(\min\{a_e,2\}),
  \label{eq:conformal-map}
\end{equation}
with $Q_e$ produced by a pure-PyTorch Householder parameterization. Cayley maps are also
implemented, while Euler rotations are restricted to $d\in\{2,3\}$. A matrix-exponential option
is rejected during configuration validation because its full-edge CUDA backward path was unsafe in
the target environment. We call the maps \emph{scaled orthogonal}, not orthogonal.

The learned first coboundary is
\begin{equation}
  (\delta_0 h^{(0)})_e=F_e^{\mathrm{tgt}}h_v^{(0)}-
  F_e^{\mathrm{src}}h_u^{(0)}.
  \label{eq:delta0}
\end{equation}
It is used as an edge feature, but not as the initial edge signal itself.

\subsection{Independent edge-cochain initialization}

The edge initializer preserves information that would be lost by forcing every initial edge state
to be exact:
\begin{equation}
  s_e^{(0)}=\operatorname{MLP}_{\mathrm{init}}\!\left(
  h_u^{(0)}\concat h_v^{(0)}\concat z_r\concat(\delta_0h^{(0)})_e
  \right).
  \label{eq:edge-init}
\end{equation}
The final block is optional through \texttt{use\_delta0\_init}. Initializing
$s^{(0)}=\delta_0h^{(0)}$ would restrict the edge signal to the image of $\delta_0$ and remove an
independent component before diffusion begins.

\subsection{Typed path instances and edge-to-face maps}

In manual mode, dataset-specific templates determine which paths are enumerated. In automatic
mode, Section~\ref{sec:ampss} constructs a fixed candidate pool. For a concrete path instance
$f=(v_0,e_1,v_1,\ldots,e_\ell,v_\ell)$, the sequence module forms tokens
\begin{equation}
  g_i=\operatorname{MLP}_{g}\!
  \left(h_{v_{i-1}}^{(0)}\concat h_{v_i}^{(0)}\concat z_{\psi(e_i)}\right),
  \qquad (o_1,\ldots,o_\ell)=\operatorname{LSTM}(g_1,\ldots,g_\ell).
\end{equation}
A learned vector $a$ gives position attention
\begin{equation}
  \pi_{fi}=\operatorname{softmax}_{i}\left(\frac{o_i^{\transpose}a}{\sqrt{d_h}}\right),
  \qquad
  A_{e_i\triangleright f}=\pi_{fi}\operatorname{reshape}(W_Ao_i+b_A,d,d).
  \label{eq:face-map}
\end{equation}
The attention weight multiplies the map itself. Maps are produced in bounded face chunks, cached
once per forward pass, and reused for all diffusion steps.

\subsection{Closure convention and the chain-condition caveat}
\label{sec:closure}

Every open path has an endpoint pair $(v_\ell,v_0)$ recorded as a synthetic closure edge. The
implementation does not add a stateful closure edge to $s$ and does not include that edge in
$\delta_1$. Instead, it adds the closed-form contribution of a frozen $\gamma I_d$ restriction to
the node Laplacian,
\begin{align}
  C_{\mathrm{cl}}(x)_{v_\ell} &\mathrel{+}=\gamma^2(x_{v_\ell}-x_{v_0}), &
  C_{\mathrm{cl}}(x)_{v_0} &\mathrel{-}=\gamma^2(x_{v_\ell}-x_{v_0}),
  \label{eq:closure}
\end{align}
for every recorded closure occurrence. The default $\gamma=10^{-6}$ yields an effective
$10^{-12}$ coupling.

Because the second operator sees only the stored path edges, independently learned restrictions do
not guarantee $\delta_1\delta_0=0$. We therefore describe the implementation as a generalized,
path-supported sheaf--Hodge system and penalize the composition during training. A future strict
cellular version would materialize a compatible closed boundary and enforce the chain condition by
construction.

\subsection{Matrix-free coboundaries and Hodge operators}

Let $\eta(e,f)\in\{-1,+1\}$ be the stored incidence sign. The second operator is
\begin{equation}
  (\delta_1s)_f=\sum_{e\in\partial_{\mathrm{path}} f}
  \eta(e,f)A_{e\triangleright f}s_e.
  \label{eq:delta1}
\end{equation}
The code supports arbitrary signs and tests a negative-incidence case. Production enumeration
follows stored forward relations, so it currently emits $+1$ for every incidence. The node and
edge operators are
\begin{align}
  L_0x &= \delta_0^{\transpose}\delta_0x+C_{\mathrm{cl}}(x),
  \label{eq:l0}\\
  L_1s &= \underbrace{\delta_0\delta_0^{\transpose}s}_{L_1^{\mathrm{down}}s}
         +\underbrace{\delta_1^{\transpose}\delta_1s}_{L_1^{\mathrm{up}}s}.
  \label{eq:l1}
\end{align}
No global $Nd\times Nd$ or $Md\times Md$ matrix is formed. Batched matrix--vector products and
scatter-add operations implement each operator and its adjoint.

\subsection{Hodge diffusion}

The model performs $K$ explicit Euler steps,
\begin{equation}
  x^{(k+1)}=x^{(k)}-\eta_0L_0x^{(k)},\qquad
  s^{(k+1)}=s^{(k)}-\eta_1L_1s^{(k)},\quad k=0,\ldots,K-1.
  \label{eq:diffusion}
\end{equation}
Learnable step sizes are softplus-parameterized and therefore positive. Optional layer dropout is
applied to both ranks after each step. Layer normalization produces $x^\star$ and $s^\star$.

The ranks are evolved in parallel under operators inferred from the same initial node state; the
implemented recurrence does not feed $s^{(k)}$ into the node update or $x^{(k)}$ into the edge
update. Higher-order information therefore reaches a task only if its readout consumes $s^\star$
or a descriptor derived from it. This observation is central to Section~\ref{sec:experiments}.

\subsection{Task readouts}
\label{sec:readouts}

\paragraph{Node classification.}
Each node receives the mean of its incident edge states,
$\bar{s}_v=|E(v)|^{-1}\sum_{e\ni v}s_e^\star$. Target-type logits are
\begin{equation}
  \widehat{y}_v=\operatorname{MLP}_{\mathrm{node}}
  (x_v^\star\concat h_v^{(0)}\concat\bar{s}_v).
  \label{eq:node-readout}
\end{equation}
This is the present readout through which the face pathway directly affects supervised predictions.

\paragraph{Link prediction.}
The implemented dot, DistMult, and concatenation-MLP decoders score pairs using $x^\star$.
DistMult uses one relation matrix per pooled target relation. Since these decoders do not consume
$s^\star$, the current link-prediction task loss does not directly supervise the upper edge
operator; link experiments validate the shared encoder and data protocol, not the central
higher-order readout claim.

\paragraph{Graph classification.}
The graph-level interface supports mean pooling of $x^\star$ and terminal heads that inspect
$x^\star$, $h^{(0)}$, and incident-edge descriptors derived from $s^\star$. The completed graph
comparison evaluates this terminal-readout setting. The planned hierarchical \thspool{} method is
a separate extension and is not part of that result.

\subsection{Training objective}

For single-label tasks, $\cL_{\mathrm{task}}$ is cross-entropy; IMDB uses multi-label binary
cross-entropy with logits. Link prediction uses binary cross-entropy. Four geometric terms are
shared across trainers:
\begin{align}
  \cL_{0} &= \lambda_0\operatorname{mean}\norm{\delta_0x^\star}_2^2, &
  \cL_{1} &= \lambda_1\operatorname{mean}\norm{\delta_1s^\star}_2^2,\\
  \cL_{\mathrm{cons}} &= \lambda_c\operatorname{mean}\norm{s^\star-\delta_0x^\star}_2^2, &
  \cL_{\mathrm{chain}} &= \lambda_{co}\operatorname{mean}
  \norm{\delta_1\delta_0h^{(0)}}_2^2.
  \label{eq:regularizers}
\end{align}
The last term is a soft chain-consistency penalty, not a proof that the learned maps form an exact
complex. The complete loss is
\begin{equation}
  \cL=\cL_{\mathrm{task}}+\cL_0+\cL_1+\cL_{\mathrm{cons}}+
  \cL_{\mathrm{chain}}+\rho(t)\cL_{\mathrm{disc}},
  \label{eq:loss}
\end{equation}
where $\rho(t)$ is zero during discovery warm-up and optionally increases linearly afterward.

\section{Adaptive Schema-to-Cell Selection}
\label{sec:ampss}

\subsection{Candidate generation and statistics}

\ampss{} begins with a directed type graph whose arcs are the stored relation tuples. Bounded
depth-first search enumerates schemas of length $\ell\in[\ell_{\min},\ell_{\max}]$, optionally
requiring the target type at both ends and allowing repeated types. Candidate count, instances per
anchor, instances per schema, and total materialized instances are separately capped. This is a
preprocessing operation and is non-differentiable.

For node classification, labels on the training mask are propagated for a fixed number of steps to
give distributions $\widehat{y}_v$. A sampled instance receives
\begin{equation}
  d(f)=\norm{\widehat{y}_{v_0}-\widehat{y}_{v_\ell}}_2^2,\qquad
  D(\sigma)=\frac{1}{|\mathcal{I}_\sigma|}\sum_{f\in\mathcal{I}_\sigma}d(f).
  \label{eq:discrepancy}
\end{equation}
Only training labels contribute to $\widehat{y}$. Link prediction uses an analogous two-channel
role signal constructed from training-positive source and target endpoints. Validation and test
positives are excluded. For graph classification, the pool is learned from training graphs only:
schema presence is counted by class and scored by the variance of class-normalized occurrence,
with a small log-frequency tie breaker. This graph-level score is task-specific and should not be
conflated with Equation~\ref{eq:discrepancy}.

The node/link path also records instance count, start and end coverage, relation diversity, type
diversity, average branching, estimated cost, and sampled endpoint sets. The structural instance
feature tensor is stored for analysis but is not currently an input to $\beta_f$.

\subsection{Greedy pool and differentiable refinement}

For node and link tasks, if $S$ is the set already selected, the score is
\begin{equation}
  q(\sigma\mid S)=
  \frac{D(\sigma)^{a_d}\,C(\sigma)^{a_c}\,V(\sigma)^{a_v}}
  {1+\max_{\sigma'\in S}J(\sigma,\sigma')},
  \label{eq:topology-score}
\end{equation}
where $C$ is end coverage, $V$ is type diversity, and $J$ is the Jaccard overlap of sampled endpoint
sets. The implementation uses $(a_d,a_c,a_v)=(1,0.5,0.3)$. After selecting $k$ schemas, scores
are recomputed leave-one-out for gate initialization.

The selected pool receives one gate per schema and one state-dependent gate per face:
\begin{align}
  \alpha_\sigma &= \operatorname{sigmoid}(\theta_\sigma),\\
  \beta_f &= \operatorname{sigmoid}\!\left(
  \operatorname{MLP}_{\beta}
  (h_{v_0}^{(0)}\concat h_{v_\ell}^{(0)}\concat e_\sigma)
  \right),\qquad
  w_f=\alpha_{\sigma(f)}\beta_f.
  \label{eq:gates}
\end{align}
The logits $\theta$ are initialized from standardized log topology scores, scaled by a temperature
and clamped. Directly setting $\theta=\log q$ placed every gate in the near-zero saturated region
in an earlier implementation.

During training, partial-schema propagation independently retains each schema with probability
$p$ and rescales its gate by $1/p$. This adapts the stabilization principle of partial message
propagation \citep{jiang2022pmmm}. The weighted operator scales each face row of $\delta_1$ by
$\sqrt{w_f}$:
\begin{equation}
  L_1^{\mathrm{up}}s=
  \sum_{\sigma}\sum_{f\in\mathcal{I}_\sigma}
  w_f B_f^{\transpose}B_f s,\qquad
  B_f s=\sum_{e\in\partial_{\mathrm{path}}f}\eta(e,f)A_{e\triangleright f}s_e.
  \label{eq:weighted-up}
\end{equation}

The discovery regularizer is
\begin{equation}
  \begin{aligned}
  \cL_{\mathrm{disc}}={}&
  \lambda_s\cL_{\mathrm{sparse\text{-}schema}}+
  \lambda_i\cL_{\mathrm{sparse\text{-}instance}}+
  \lambda_r\cL_{\mathrm{redundancy}}\\
  &+\lambda_v\cL_{\mathrm{coverage}}+
  \lambda_{\mathrm{cost}}\cL_{\mathrm{cost}}+
  \lambda_H\cL_{\mathrm{entropy}}.
  \end{aligned}
  \label{eq:discovery-loss}
\end{equation}
Schema and instance sparsity minimize mean gate values; coverage is a negative reward; cost
penalizes long, frequent schemas; and the positive entropy coefficient sharpens schema gates. The
inexpensive redundancy term uses a coverage profile rather than an exact pairwise instance-overlap
matrix. Instance sparsity consumes a gradient-attached gate mean; a detached copy is logging-only.

\begin{algorithm}[t]
  \caption{\ampss{} preprocessing and refinement}
  \label{alg:ampss}
  \begin{algorithmic}[1]
    \Require Relation tuples, training split, candidate and instance budgets
    \State Build the directed type graph and enumerate bounded typed schemas
    \State Sample instances and compute task-specific training-only statistics
    \State Greedily retain a compact, non-redundant schema pool
    \State Materialize path groups, incidence signs, and closure endpoint pairs
    \State Initialize $\alpha$ logits from standardized log topology scores
    \For{each training epoch}
      \State Use uniform faces during warm-up; otherwise compute $\alpha$ and $\beta$
      \State Apply partial-schema propagation and set $w_f=\alpha_{\sigma(f)}\beta_f$
      \State Run matrix-free Hodge diffusion and the task readout
      \State Optimize Equation~\ref{eq:loss}
    \EndFor
  \end{algorithmic}
\end{algorithm}

\section{Operator Properties and Computational Cost}
\label{sec:analysis}

\begin{proposition}[Positive semidefiniteness]
For fixed restriction maps and nonnegative face weights, $L_0$ and $L_1$ in
Equations~\ref{eq:l0}--\ref{eq:weighted-up} are self-adjoint positive semidefinite.
\end{proposition}
\begin{proof}
For every node cochain $x$,
$\langle x,L_0x\rangle=\norm{\delta_0x}_2^2+
\sum_c\gamma^2\norm{x_{t(c)}-x_{s(c)}}_2^2\geq0$.
For every edge cochain $s$,
$\langle s,L_1s\rangle=\norm{\delta_0^{\transpose}s}_2^2+
\sum_f w_f\norm{B_f s}_2^2\geq0$. Each expression is a sum of Gram terms.
\end{proof}

The sigmoid parameterization makes $w_f\geq0$. Partial-schema propagation preserves
nonnegativity, even when rescaling makes a retained weight larger than one.

\begin{proposition}[One-step energy stability]
For a fixed positive-semidefinite operator $L$, $z^+=z-\eta Lz$ is non-expansive in Euclidean norm
when $0\leq\eta\leq2/\lambda_{\max}(L)$.
\end{proposition}
\begin{proof}
Diagonalize $L=U\Lambda U^{\transpose}$. Each spectral component is multiplied by
$1-\eta\lambda_i$, whose magnitude is at most one under the stated condition.
\end{proof}

The implementation enforces $\eta>0$ but not the spectral upper bound. The proposition is a
sufficient condition, not an automatic guarantee of every trained configuration. Layer
normalization and gradient clipping are practical safeguards.

\paragraph{Permutation behaviour.}
Predictors are shared by type or relation, and adjoints use indexed scatter-add. A consistent
relabeling of nodes, edges, and path instances therefore permutes intermediate cochains
equivariantly. Target selection and symmetric pooling remove the final indexing dependence.

\paragraph{Cost.}
Once dense $d\times d$ restrictions are available, one $L_0$ application costs $O(Md^2)$ and one
$L_1$ application costs $O((M+I)d^2)$. A $K$-step rollout costs $O(K(M+I)d^2)$ for operator
application. Dense map storage is $O(Md^2+Id^2)$, and Householder/Cayley construction can dominate
at larger $d$. The code chunks face maps and rejects CUDA configurations whose estimated
node-to-edge materialization exceeds a configurable fraction (default $0.20$) of device memory.

\section{Experimental Evaluation}
\label{sec:experiments}

The evaluation follows the breadth of HetSheaf \citep{braithwaite2024hetsheaf} while adding
clustering and a direct test of automatic schema-to-cell discovery. All values in
Sections~\ref{sec:results}--\ref{sec:analysis-results} are percentages and reproduce the completed
averages recorded by the project experiment reports. The headline \method{} row uses the full
\ampss{} configuration, with both schema gates $\alpha$ and instance gates $\beta$ active.

\subsection{Research questions}
\begin{enumerate}
  \item[RQ1.] Does path-supported sheaf--Hodge diffusion improve node classification across
  heterogeneous homophily regimes?
  \item[RQ2.] Does the same representation transfer to graph classification, node clustering, and
  heterogeneous link prediction?
  \item[RQ3.] Does \ampss{} improve over both removing meta-path cells and prescribing them by hand,
  and does instance-level gating add value beyond schema selection?
  \item[RQ4.] Does the higher-order construction remain competitive in runtime, parameter count, and
  peak memory?
\end{enumerate}

\subsection{Datasets and protocols}
\begin{table}[t]
  \centering
  \caption{Datasets, sources, and evaluation measures.}
  \label{tab:datasets}
  \small
  \begin{tabularx}{\textwidth}{llXX}
    \toprule
    \textbf{Task} & \textbf{Datasets} & \textbf{Source} & \textbf{Measures} \\
    \midrule
    Node classification & ACM, DBLP, IMDB & HGB & Macro-F1, Micro-F1 \\
    Graph classification & MUTAG, PROTEINS, ENZYMES, NCI1 & TUDataset & Macro-F1, Micro-F1 \\
    Node clustering & ACM, Freebase, AMiner, IMDB & Heterogeneous benchmarks & NMI, ARI \\
    Link prediction & LastFM, PubMed, Amazon, YouTube & HGB-format benchmark files & AUPR, AUROC \\
    \bottomrule
  \end{tabularx}
\end{table}

For HGB node classification, all graph relations remain available in the transductive graph while
schema discovery sees training labels only. IMDB retains its multi-label genre targets and uses
sigmoid predictions. The node results average the recorded repeated runs on the official target
masks. Graph classification follows the recorded ten-fold protocol; types are recovered from TU
node and edge labels when available, otherwise the graph uses the trivial one-type annotation. The
clustering study evaluates learned target-node representations with NMI and ARI. Link prediction
uses the official target relations and fixed test triples supplied with the benchmark files; held-out
target edges are excluded from message passing, and negatives are fixed for evaluation. Appendix
\ref{app:datasets} records the precise target relations and schema presets.

The graph-classification results below concern the completed graph-level \method{} evaluation. They
do not evaluate \thspool{}, which is a separate hierarchical coarsening proposal retained only as
the placeholder in Appendix~\ref{app:thspool}.

\subsection{Baselines, ablations, and model selection}

The node and graph comparisons include GCN/GAT \citep{kipf2017gcn,velickovic2018gat}, HAN/HGT
\citep{wang2019han,hu2020hgt}, Neural Sheaf Diffusion \citep{bodnar2022neural}, HetSheaf
\citep{braithwaite2024hetsheaf}, and heterophily-aware methods
\citep{guo2023hdhgr,lu2024h2sgnn,shen2024latgrl}. The link study adds GIN, GraphSAGE, and R-GCN
\citep{xu2019gin,hamilton2017sage,schlichtkrull2018rgcn}. Clustering compares with HGMAE, RMR,
DMGI, HDMI, HeCo/HeCo++, HERO, CTC, and the two HGMS variants \citep{wang2025hgms}.
Automatic-selection context includes LMSPS and PMMM \citep{yan2024lmsps,jiang2022pmmm}.

The primary discovery ablation compares no meta-path cells, predefined schemas, \ampss{} with
$\alpha$ alone, and full \ampss{} with $\alpha\beta$. Additional implementation switches remove
faces, suppress $\delta_1$, replace learned node-to-edge maps with identity maps, or set all gates
to one; these controls are described in Appendix~\ref{app:ablations}.

All experiments use validation-selected checkpoints, Adam or AdamW, gradient clipping, and cosine
or reduce-on-plateau scheduling. W\&B tuning executes each trial in a fresh worker, fingerprints
resume state, and rejects stale search spaces. Every run saves its resolved configuration, log,
JSONL metrics, last checkpoint, and best checkpoint.

\section{Results}
\label{sec:results}

\subsection{Node classification}

Table~\ref{tab:node-cls} compares average Macro-F1 and Micro-F1 on DBLP, ACM, and IMDB. The
parenthesized values in the dataset headers are the dataset-level homophily indicators recorded by
the evaluation report. \method{} ranks first in every column.

\begin{table}[H]
  \centering
  \caption{Node classification: average Macro-F1 and Micro-F1. Best values are bold.}
  \label{tab:node-cls}
  \scriptsize
  \resizebox{\textwidth}{!}{%
  \begin{tabular}{llcccccc}
    \toprule
    & & \multicolumn{2}{c}{DBLP (59.76)} & \multicolumn{2}{c}{ACM (72.39)} & \multicolumn{2}{c}{IMDB (50.04)} \\
    \cmidrule(lr){3-4}\cmidrule(lr){5-6}\cmidrule(lr){7-8}
    \textbf{Family} & \textbf{Model} & Macro & Micro & Macro & Micro & Macro & Micro \\
    \midrule
    Homogeneous & GCN & 90.77 & 90.28 & 91.63 & 91.38 & 62.33 & 67.69 \\
    Homogeneous & GAT & 94.39 & 92.52 & 92.07 & 90.71 & 62.13 & 68.09 \\
    Sheaf/spectral & NSD & 92.96 & 93.42 & 92.33 & 92.05 & 61.61 & 65.23 \\
    Sheaf/spectral & CSNN & 93.73 & 92.12 & 93.13 & 92.81 & 62.92 & 62.59 \\
    Sheaf/spectral & ChebNetII & 92.89 & 92.67 & 91.70 & 91.59 & 62.97 & 68.47 \\
    Heterogeneous & HAN & 91.78 & 92.24 & 90.88 & 91.61 & 61.95 & 67.76 \\
    Heterogeneous & SeHGNN & 95.52 & 95.48 & 94.38 & 94.04 & 72.01 & 72.59 \\
    Heterogeneous & HiNormer & 94.06 & 94.14 & 91.87 & 89.64 & 64.13 & 67.53 \\
    Heterogeneous & PSHGCN & 95.58 & 95.70 & 93.66 & 94.70 & 71.82 & 74.34 \\
    Heterophily-aware & HDHGR & 95.35 & 95.00 & 93.99 & 94.13 & 59.41 & 59.63 \\
    Heterophily-aware & Hetero2Net & 93.54 & 93.54 & 94.05 & 92.87 & 65.34 & 70.63 \\
    Heterophily-aware & H2SGNN & 95.82 & 95.11 & 94.90 & 94.08 & 74.23 & 75.29 \\
    Heterogeneous sheaf & HetSheaf & 93.87 & 93.77 & 95.09 & 94.94 & 64.65 & 67.18 \\
    Heterophily-aware & LatGRL & 94.32 & 94.51 & 93.47 & 92.29 & 58.27 & 58.81 \\
    \midrule
    Ours & \textbf{\method{} + \ampss{}} & \textbf{96.01} & \textbf{96.31} & \textbf{95.91} & \textbf{95.72} & \textbf{75.05} & \textbf{76.68} \\
    \bottomrule
  \end{tabular}%
  }
\end{table}

On DBLP, \method{} exceeds the strongest competing Macro-F1 by $0.19$ points and Micro-F1 by
$0.61$ points. The corresponding margins are $0.82/0.78$ on ACM and $0.82/1.39$ on IMDB. The IMDB
gain is particularly informative because its reported homophily indicator is the lowest of the
three and its targets are multi-label. The result is consistent with the intended role of learned
restrictions: comparison occurs after relation-conditioned transport, while path instances preserve
longer typed context in the upper edge operator.

\subsection{Graph classification}

Table~\ref{tab:graph-cls} reports ten-fold graph-classification averages. \method{} reaches the
shared ceiling on MUTAG and leads on PROTEINS, ENZYMES, and NCI1. These experiments use the recorded
graph-level readout and are distinct from the future \thspool{} hierarchy.

\begin{table}[H]
  \centering
  \caption{Graph classification: average Macro-F1 and Micro-F1. Best values are bold.}
  \label{tab:graph-cls}
  \scriptsize
  \resizebox{\textwidth}{!}{%
  \begin{tabular}{llcccccccc}
    \toprule
    & & \multicolumn{2}{c}{MUTAG} & \multicolumn{2}{c}{PROTEINS} & \multicolumn{2}{c}{ENZYMES} & \multicolumn{2}{c}{NCI1} \\
    \cmidrule(lr){3-4}\cmidrule(lr){5-6}\cmidrule(lr){7-8}\cmidrule(lr){9-10}
    \textbf{Family} & \textbf{Model} & Ma. & Mi. & Ma. & Mi. & Ma. & Mi. & Ma. & Mi. \\
    \midrule
    Homogeneous & GCN & 89.58 & 90.00 & 75.92 & 79.83 & 12.75 & 28.33 & 64.17 & 64.96 \\
    Homogeneous & GAT & 89.58 & 90.00 & 75.27 & 78.38 & 14.03 & 16.55 & 64.17 & 64.96 \\
    Sheaf & NSD & \textbf{100} & \textbf{100} & 73.60 & 75.89 & 33.65 & 35.00 & 72.30 & 72.30 \\
    Sheaf & CSNN & \textbf{100} & \textbf{100} & 76.64 & 77.92 & 35.15 & 37.22 & 74.11 & 74.30 \\
    Heterogeneous & HAN & 82.91 & 85.00 & 79.19 & 82.22 & 46.24 & 45.02 & 77.60 & 77.50 \\
    Heterogeneous & PSHGCN & 94.88 & 95.00 & 70.61 & 72.32 & 20.18 & 28.33 & 64.17 & 64.96 \\
    Heterophily-aware & HDHGR & 84.65 & 85.00 & 77.64 & 81.25 & 48.27 & 48.33 & 79.29 & 79.32 \\
    Heterophily-aware & Hetero2Net & 89.58 & 90.00 & 74.66 & 78.57 & 8.53 & 18.33 & 63.67 & 65.45 \\
    Heterophily-aware & H2SGNN & 94.88 & 95.00 & 70.61 & 72.32 & 20.18 & 28.33 & 64.17 & 64.96 \\
    Heterogeneous sheaf & HetSheaf & \textbf{100} & \textbf{100} & 80.56 & 82.14 & 62.35 & 63.33 & 79.46 & 79.46 \\
    Heterophily-aware & LatGRL & 90.88 & 91.01 & 71.61 & 72.32 & 25.41 & 29.03 & 77.98 & 78.10 \\
    \midrule
    Ours & \textbf{\method{} + \ampss{}} & \textbf{100} & \textbf{100} & \textbf{84.44} & \textbf{86.23} & \textbf{64.52} & \textbf{65.10} & \textbf{81.22} & \textbf{81.52} \\
    \bottomrule
  \end{tabular}%
  }
\end{table}

Relative to the next-best HetSheaf row, the improvement is $3.88/4.09$ points on PROTEINS,
$2.17/1.77$ on ENZYMES, and $1.76/2.06$ on NCI1 for Macro-/Micro-F1. The larger gains on PROTEINS
and ENZYMES are consistent with a benefit from representations of distributed graph motifs. Since
the table compares complete models, causal attribution to a particular rank is left to component
ablations rather than inferred from the headline scores alone.

\subsection{Node clustering}

Table~\ref{tab:clustering} evaluates the learned target-node representation using normalized mutual
information (NMI) and adjusted Rand index (ARI). \method{} is best in every column.

\begin{table}[H]
  \centering
  \caption{Node clustering: NMI and ARI. Best values are bold.}
  \label{tab:clustering}
  \scriptsize
  \resizebox{\textwidth}{!}{%
  \begin{tabular}{lcccccccc}
    \toprule
    & \multicolumn{2}{c}{ACM} & \multicolumn{2}{c}{Freebase} & \multicolumn{2}{c}{AMiner} & \multicolumn{2}{c}{IMDB} \\
    \cmidrule(lr){2-3}\cmidrule(lr){4-5}\cmidrule(lr){6-7}\cmidrule(lr){8-9}
    \textbf{Model} & NMI & ARI & NMI & ARI & NMI & ARI & NMI & ARI \\
    \midrule
    HGMAE & 64.57 & 67.85 & 17.18 & 17.94 & 34.70 & 34.38 & 2.76 & 2.62 \\
    RMR & 54.37 & 45.65 & 18.33 & 18.15 & 30.22 & 31.59 & 2.59 & 2.47 \\
    DMGI & 56.49 & 51.06 & 16.41 & 16.13 & 19.24 & 20.09 & 2.15 & 1.51 \\
    HDMI & 57.91 & 54.75 & 16.53 & 17.47 & 17.56 & 15.04 & 2.07 & 1.79 \\
    HeCo & 56.87 & 56.94 & 16.14 & 16.90 & 32.26 & 28.64 & 2.72 & 2.44 \\
    HeCo++ & 60.82 & 60.09 & 16.67 & 17.48 & 38.07 & 36.44 & 2.92 & 2.58 \\
    HERO & 58.41 & 57.89 & 15.34 & 15.56 & 30.11 & 31.48 & 2.19 & 2.05 \\
    CTC & 63.26 & 69.30 & 18.70 & 21.31 & 35.21 & 34.63 & 2.86 & 2.55 \\
    HGMS-N & 72.44 & 75.53 & 21.35 & 22.49 & 42.87 & 43.67 & 6.66 & 5.25 \\
    HGMS-C & 71.97 & 74.18 & 22.28 & 23.14 & 46.56 & 50.28 & 6.61 & 6.77 \\
    \midrule
    \textbf{\method{} + \ampss{}} & \textbf{80.69} & \textbf{79.50} & \textbf{24.01} & \textbf{24.81} & \textbf{52.32} & \textbf{56.20} & \textbf{7.59} & \textbf{7.17} \\
    \bottomrule
  \end{tabular}%
  }
\end{table}

The gains over the best competing entry are $8.25/3.97$ points on ACM, $1.73/1.67$ on Freebase,
$5.76/5.92$ on AMiner, and $0.93/0.40$ on IMDB. Unlike classification, this comparison does not let
the final linear decision boundary absorb representational defects. The consistent ranking therefore
supports the stronger interpretation that the learned sheaf--Hodge geometry improves the embedding
space itself.

\subsection{Heterogeneous link prediction}

Table~\ref{tab:link-pred} reports AUPR and AUROC. \method{} gives the strongest result in all eight
columns, including exact $100$ on LastFM and non-saturated gains on the three other benchmarks.

\begin{table}[H]
  \centering
  \caption{Heterogeneous link prediction: AUPR and AUROC. Best values are bold.}
  \label{tab:link-pred}
  \scriptsize
  \resizebox{\textwidth}{!}{%
  \begin{tabular}{llcccccccc}
    \toprule
    & & \multicolumn{2}{c}{LastFM} & \multicolumn{2}{c}{PubMed} & \multicolumn{2}{c}{Amazon} & \multicolumn{2}{c}{YouTube} \\
    \cmidrule(lr){3-4}\cmidrule(lr){5-6}\cmidrule(lr){7-8}\cmidrule(lr){9-10}
    \textbf{Family} & \textbf{Model} & AUPR & AUROC & AUPR & AUROC & AUPR & AUROC & AUPR & AUROC \\
    \midrule
    Homogeneous & GCN & 97.06 & 96.96 & 96.05 & 92.93 & 92.40 & 89.66 & 89.30 & 83.26 \\
    Homogeneous & GAT & 63.06 & 51.14 & 96.07 & 93.27 & 88.62 & 86.10 & 66.76 & 58.40 \\
    Homogeneous & GIN & 97.99 & 97.75 & 96.14 & 92.87 & 92.10 & 89.62 & 88.64 & 82.97 \\
    Homogeneous & GraphSAGE & 95.11 & 92.72 & 94.50 & 89.79 & 90.79 & 87.95 & 84.54 & 80.20 \\
    Heterogeneous & R-GCN & 97.89 & 97.93 & 95.03 & 93.15 & 93.37 & 91.42 & 91.68 & 88.33 \\
    Heterogeneous & HAN & 88.26 & 83.32 & 95.78 & 93.62 & 95.24 & 93.64 & 93.46 & 90.10 \\
    Heterogeneous & HGT & 86.02 & 74.55 & 97.58 & 96.48 & 96.15 & 95.66 & 94.95 & 92.51 \\
    Sheaf & Sheaf-NSD & 97.16 & 96.58 & 96.41 & 93.65 & 92.63 & 90.04 & 89.32 & 83.36 \\
    Sheaf & HetSheaf & 98.51 & 98.33 & 96.36 & 93.98 & 93.76 & 91.71 & 91.27 & 87.33 \\
    \midrule
    Ours & \textbf{\method{} + \ampss{}} & \textbf{100} & \textbf{100} & \textbf{99.63} & \textbf{99.28} & \textbf{99.71} & \textbf{98.34} & \textbf{98.37} & \textbf{96.30} \\
    \bottomrule
  \end{tabular}%
  }
\end{table}

The margins over the strongest alternative are $1.49/1.67$ points on LastFM, $2.05/2.80$ on
PubMed, $3.56/2.68$ on Amazon, and $3.42/3.79$ on YouTube. The fact that both ranking-sensitive AUPR
and threshold-free AUROC improve argues against a result driven only by a favorable operating
threshold. Because the implemented link decoder scores $x^\star$ rather than $s^\star$, these
columns establish the transfer of the shared typed sheaf encoder; they are not presented as a
causal ablation of the upper edge term.

\section{Analysis}
\label{sec:analysis-results}

\subsection{Contribution of automatic schema-to-cell selection}

Table~\ref{tab:ablation-schema} isolates the discovery mechanism on DBLP. Removing meta-path cells
reduces Macro-/Micro-F1 by $2.16/2.59$ points relative to full \ampss{}. Hand-specified schemas
recover most of this gap. Replacing them with automatic schema selection and an $\alpha$ gate adds
$0.13$ Macro-F1 but changes Micro-F1 by $-0.09$; adding the instance gate $\beta$ then contributes a
further $0.13/0.39$ points and produces the best row. This non-monotonic intermediate result is
useful: automatic schema ranking alone is not sufficient, whereas instance-level refinement can
suppress uninformative realizations of an otherwise useful type schema.

\begin{table}[H]
  \centering
  \caption{Schema-to-cell ablation on DBLP node classification.}
  \label{tab:ablation-schema}
  \small
  \begin{tabular}{lcc}
    \toprule
    \textbf{Configuration} & Macro-F1 & Micro-F1 \\
    \midrule
    No meta-path cells & 93.85 & 93.72 \\
    Predefined meta-paths & 95.75 & 96.01 \\
    \ampss{}, $\alpha$ only & 95.88 & 95.92 \\
    \textbf{Full \ampss{}, $\alpha+\beta$} & \textbf{96.01} & \textbf{96.31} \\
    \bottomrule
  \end{tabular}
\end{table}

\subsection{Efficiency}

The matrix-free formulation does not materialize global block Laplacians; its practical resource
profile is summarized in Table~\ref{tab:efficiency}. On ACM, \method{} is $16.8\%$ faster than the
next-fastest R-GCN, uses $13.5\%$ fewer parameters than the next-smallest Sheaf-NSD, and consumes
$17.3\%$ less peak memory than that model. Relative to HGT, it is $55.8\times$ faster, has
$53.7\times$ fewer parameters, and uses $7.6\times$ less memory in this comparison.

\begin{table}[H]
  \centering
  \caption{Efficiency comparison on ACM. Best values are bold.}
  \label{tab:efficiency}
  \small
  \begin{tabular}{lccc}
    \toprule
    \textbf{Model} & Runtime (s) & Parameters & Memory (MB) \\
    \midrule
    R-GCN & 22.1 & 43M & 1{,}240 \\
    HAN & 28.9 & 2.2M & 1{,}780 \\
    HGT & 1{,}025.8 & 7.2M & 6{,}120 \\
    Sheaf-NSD & 47.6 & 155K & 980 \\
    \midrule
    \textbf{\method{}} & \textbf{18.39} & \textbf{134K} & \textbf{810} \\
    \bottomrule
  \end{tabular}
\end{table}

\subsection{Cross-task interpretation}

Three patterns recur. First, the advantage is not confined to a supervised node classifier:
clustering improves in all eight columns, suggesting a change in representation geometry rather
than only classifier capacity. Second, gains remain visible for relation prediction and molecular
or protein graph labels, showing that the shared typed sheaf encoder transfers across prediction
units. These heads expose the node and edge ranks differently, so their headline gains are not by
themselves causal evidence for $\delta_1$. Third, the strongest relative improvements occur on
AMiner clustering, PROTEINS graph classification, and Amazon/YouTube link prediction, tasks where
local adjacency alone is often semantically incomplete. The \ampss{} ablation supplies the direct
mechanism-level evidence: retaining higher-order supports matters, automatic schema choice helps,
and instance filtering adds a distinct final gain.

\section{Discussion}
\label{sec:discussion}

The principal conceptual step is the promotion of selected path instances from routing rules to
terms in a learned edge operator. Node-to-edge maps allow relation-specific comparison of
heterogeneous endpoints; edge-to-face maps preserve the ordered typed boundary; and nonnegative
gates select which higher-order supports shape diffusion. The empirical picture is consistent with
this division of labor. The broad benchmark gains show transfer across prediction units, while the
DBLP ablation shows that both the existence of higher-order supports and their hierarchical
selection contribute beyond a no-meta-path construction.

The results also clarify the relationship to adjacent methods. Against conventional HGNNs, the
gain comes without reducing a selected path to a disposable aggregation channel. Against Neural
Sheaf Diffusion and HetSheaf, the model adds an independently represented edge rank and an upper
Hodge term. Against schema-search methods, the selected object is not merely an architecture choice:
it becomes a weighted operator summand. This difference is especially visible in the clustering
results, where no supervised classifier can compensate for a poorly structured embedding.

Several limitations concern the mathematical form rather than experimental completeness. Closure
pairs affect $L_0$ but are absent from $\delta_1$, so the chain condition is softly regularized rather
than guaranteed structurally. Candidate preselection is deliberately non-differentiable, followed
by differentiable refinement over a fixed pool. Signed incidence is implemented, yet the production
DFS traverses stored relations forward and therefore emits positive signs. Finally, every active
incidence may carry a dense $d\times d$ map, making face budgets and stalk dimension the principal
scaling variables even though the ACM resource measurements are favorable.

The only method component intentionally left incomplete is \thspool{} (internally S$^2$Pool-H2),
the proposed hierarchical sheaf pooling extension. The graph-classification results in
Table~\ref{tab:graph-cls} use the completed graph-level evaluation path and do not constitute an
evaluation of that hierarchy. Immediate future work is therefore concentrated on \thspool{}, a
strict closed-boundary construction with a structural chain condition, and lower-rank restriction
parameterizations. Calibrated hard-concrete gates may also turn the continuous \ampss{} selection
into a more discrete and interpretable complex.

\section{Conclusion}

We presented \method, a heterogeneous sheaf--Hodge architecture that retains selected typed path
instances as higher-order supports for an explicit edge diffusion operator. Learned
relation-conditioned node-to-edge maps, sequence-conditioned edge-to-face maps, and matrix-free
Hodge terms define the backbone. \ampss{} complements it with training-only schema preselection
and differentiable schema/instance refinement inside the upper edge Laplacian. Across node and
graph classification, node clustering, and heterogeneous link prediction, the complete evaluation
places \method{} first or tied for first in every reported comparison column, while the DBLP
ablation isolates gains from both automatic schema selection and instance gating. The ACM resource
study further shows that the higher-order construction can be realized without a parameter or
memory penalty relative to the compared methods. These results support the central claim that typed
paths are most useful when retained as operator-level higher-order structure rather than consumed
only as aggregation routes. \thspool{} remains the sole placeholder and is reserved for future work.

\bibliographystyle{plainnat}
\bibliography{references}

\appendix

\section{Notation and Full Forward Pass}
\label{app:notation}

\begin{table}[h]
  \centering
  \caption{Main symbols used in the paper.}
  \small
  \begin{tabularx}{\textwidth}{lX}
    \toprule
    \textbf{Symbol} & \textbf{Meaning} \\
    \midrule
    $G=(V,E,\phi,\psi)$ & Directed heterogeneous graph with node- and relation-type maps. \\
    $N,M,F,I$ & Numbers of nodes, stored edges, path instances, and path-edge incidences. \\
    $d,d_r,d_h$ & Stalk, relation-embedding, and sequence hidden dimensions. \\
    $h^{(0)},x^\star,s^{(0)},s^\star$ & Initial/final node and edge cochains. \\
    $F_e^{\mathrm{src}},F_e^{\mathrm{tgt}}$ & Learned node-to-edge restriction maps. \\
    $A_{e\triangleright f}$ & Learned edge-to-face restriction map. \\
    $\delta_0,\delta_1$ & First and path-supported second coboundary operators. \\
    $L_0,L_1$ & Node and edge Hodge-style Laplacians. \\
    $\alpha_\sigma,\beta_f,w_f$ & Schema, instance, and effective face weights. \\
    $\gamma$ & Frozen synthetic-closure restriction scale. \\
    \bottomrule
  \end{tabularx}
\end{table}

\begin{algorithm}[h]
  \caption{\method{} forward pass for node classification}
  \begin{algorithmic}[1]
    \Require Typed graph tensors, face groups, closure pairs, target indices, epoch
    \State Encode typed nodes using Equation~\ref{eq:node-encoder}
    \State Predict $F^{\mathrm{src}},F^{\mathrm{tgt}}$ and compute $\delta_0h^{(0)}$
    \State Initialize $s^{(0)}$ with Equation~\ref{eq:edge-init}
    \State Encode ordered path boundaries and cache $A_{e\triangleright f}$
    \State Build uniform weights, or $\alpha\beta$ weights after discovery warm-up
    \For{$k=0$ to $K-1$}
      \State Apply $L_0$, including the closure correction
      \State Apply $L_1^{\mathrm{down}}+L_1^{\mathrm{up}}$
      \State Take the Euler step and apply layer dropout
    \EndFor
    \State Layer-normalize $x^\star,s^\star$ and compute target logits
    \State Return logits, cochain diagnostics, and gate summaries
  \end{algorithmic}
\end{algorithm}

\section{Dataset and Task Details}
\label{app:datasets}

\begin{table}[h]
  \centering
  \caption{Checked-in manual node-classification schema presets.}
  \small
  \begin{tabular}{llll}
    \toprule
    \textbf{Dataset} & \textbf{Target type} & \textbf{Manual schemas} & \textbf{Labels} \\
    \midrule
    ACM & paper & PAP, PSP & single-label, 3 classes \\
    DBLP & author & APA, APTPA & single-label, 4 classes \\
    IMDB & movie & MAM, MDM & multi-label, 5 genres \\
    \bottomrule
  \end{tabular}
\end{table}

LastFM targets $(\text{user},\mathrm{rel}_0,\text{artist})$; PubMed targets
$(\text{disease},\mathrm{rel}_2,\text{disease})$; YouTube targets one user--user relation at a
time; and HGB Amazon pools its two product--product relations while passing relation IDs to
DistMult. AmazonBook is a separate user--book dataset and is never compared under the name
``Amazon'' to an HGB Amazon baseline.

For graph discovery, training-graph paths are grouped by the alternating typed signature.
Candidate presence is counted at most once per graph for class association, while total instance
count is retained as a cost/frequency feature. Selected training-derived names filter every
mini-batch. No validation or test graph label enters template construction.

\section{Hyperparameters and Reproducibility}
\label{app:hyperparameters}

\begin{table}[h]
  \centering
  \caption{Base values in checked-in ACM/DBLP/IMDB configurations. Tuning overrides are stored in
  each resolved run configuration.}
  \small
  \begin{tabular}{lll}
    \toprule
    \textbf{Group} & \textbf{Field} & \textbf{Base value} \\
    \midrule
    Stalks & $d,d_r,d_h$ & $32,16,64$ \\
    Diffusion & $K,\eta_0,\eta_1$ & $4,0.05,0.05$ \\
    Maps & family & scaled Householder \\
    Networks & edge/init/face/classifier hidden & $128$ each \\
    Geometry & $\lambda_0,\lambda_1,\lambda_c,\lambda_{co}$ & $0,10^{-3},10^{-2},10^{-2}$ \\
    Optimizer & name, learning rate, weight decay & AdamW, $10^{-3}$, $10^{-4}$ \\
    Training & epochs, gradient clip, base seed & $200,1.0,42$ \\
    Discovery & schema length, selected $k$ & $2$--$4$, dataset dependent ($6$ or $8$) \\
    Discovery & partial probability, warm-up & $0.8$, $5$ epochs in AMPSS configs \\
    \bottomrule
  \end{tabular}
\end{table}

The driver saves \texttt{config.yaml}, \texttt{run.log}, \texttt{metrics.jsonl},
\texttt{checkpoints/last.pt}, and \texttt{checkpoints/best.pt}. Validation metrics alone select the
best epoch. Tuning state is fingerprinted per seed and search space. The test suite uses
\texttt{unittest}; on the target Windows machine, PyG imports can be blocked by Application Control
at \texttt{xxhash}, an environment-policy failure rather than a model failure.

\section{Additional Diagnostic Ablations}
\label{app:ablations}

\begin{table}[h]
  \centering
  \caption{Complementary diagnostic controls supported by the implementation.}
  \small
  \begin{tabularx}{\textwidth}{lX}
    \toprule
    \textbf{Ablation} & \textbf{Interpretation} \\
    \midrule
    No faces & Reduces to node/edge sheaf diffusion without path-supported upper terms. \\
    No $\delta_1$ & Keeps data construction while deleting $L_1^{\mathrm{up}}$. \\
    Identity node-to-edge maps & Tests whether learned heterogeneous transport is necessary. \\
    Uniform $\alpha,\beta$ & Separates discovery weights from face support. \\
    Manual schemas & Compares automatic construction to expert templates. \\
    $\alpha$ only & Separates schema-level from instance-level refinement, as in Table~\ref{tab:ablation-schema}. \\
    Equal-parameter control & Tests topology rather than additional capacity. \\
    Closure sweep & Measures whether the very small default $\gamma$ has a detectable effect. \\
    Face/path budget sweep & Reports the accuracy--memory trade-off as $F$ and $I$ change. \\
    \bottomrule
  \end{tabularx}
\end{table}

\section{Implementation Traceability}
\label{app:traceability}

\begin{table}[h]
  \centering
  \caption{Source files grounding the methodological statements.}
  \small
  \begin{tabularx}{\textwidth}{lX}
    \toprule
    \textbf{Component} & \textbf{Authoritative implementation} \\
    \midrule
    Typed encoding & \texttt{src/models/node\_encoder.py} \\
    Node--edge and edge--face maps & \texttt{src/models/sheaf\_maps.py} \\
    Operators and closure correction & \texttt{src/models/coboundary\_ops.py} \\
    Euler diffusion & \texttt{src/models/diffusion.py} \\
    Schema candidates and scoring & \texttt{src/data/schema\_discovery.py} \\
    Discovery gates and losses & \path{src/models/schema_gates.py}; \path{src/models/discovery_losses.py} \\
    Task readouts & \texttt{src/models/readout.py}, \texttt{link\_readout.py}, \texttt{graph\_pooling.py} \\
    End-to-end wiring & \texttt{src/models/h2snn\_model.py} \\
    \bottomrule
  \end{tabularx}
\end{table}

\section{Status of THS-Pool}
\label{app:thspool}

\thspool{} is the working name used here for the hierarchical, typed, higher-order pooling project
documented internally as S$^2$Pool-H2. Its intended design performs type-aware node selection,
constructs a coarse typed graph, re-enumerates path supports at the coarse resolution, re-encodes
the coarse stalks, and reads node, edge, and face ranks across resolutions. Only the
\texttt{encode\_from\_stalk} backbone entry point and terminal diagnostic heads are currently in the
codebase; the pooler, coarse driver, and final controlled evaluation are not complete. This
appendix reserves the connection without treating \thspool{} as part of the present method or
results.

\end{document}
